\documentclass[acmsmall,nonacm,11pt,natbib=false]{acmart}
\geometry{letterpaper,left=1in,right=1in,top=1in,bottom=1in}
\settopmatter{printacmref=false}
\setcopyright{none}
\pagestyle{plain}
\raggedbottom
\AtBeginDocument{\let\small\normalsize\let\footnotesize\normalsize}
\usepackage{listings}
\lstset{language=Python,basicstyle=\ttfamily\normalsize,
  columns=fullflexible,keepspaces=true,showstringspaces=false,
  frame=single,framerule=0.3pt,framesep=6pt,aboveskip=12pt,
  belowskip=12pt,xleftmargin=7pt,xrightmargin=7pt,breaklines=false}
\newcommand{\wild}{\mathord{\diamond}}
\title{A Universal Partial Cycle with Three Wildcards per Window}
\author{Zijian Zeng}
\affiliation{\institution{UCSI University}\country{Malaysia}}
\email{1002266693@ucsiuniversity.edu.my}
\authorsaddresses{Zijian Zeng, UCSI University.\newline
  \href{mailto:1002266693@ucsiuniversity.edu.my}{\nolinkurl{1002266693@ucsiuniversity.edu.my}}}
\date{}

\begin{document}
\begin{abstract}
A universal partial cycle covers every word of a fixed length exactly once
by cyclic windows, allowing a wildcard to match any alphabet symbol. The
number of wildcards in each window is its diamondicity. Fillmore,
Goeckner, Kirsch, Martin and McGinnis asked whether a nontrivial universal
partial cycle of diamondicity greater than one exists. We answer this
question affirmatively by giving a binary cycle of length 512 whose
length-12 windows each contain three wildcards. Its 512 windows partition
the 4096 binary words into sets of eight completions.

The complete cycle is specified by 128 octal digits. A short, solver-free
Python program checks every pair of cyclic windows and every binary
12-word, giving a reproducible finite existence proof. The cycle also
realizes the regime in which the window length does not divide the cycle
length.
\end{abstract}
\maketitle
\pagestyle{plain}
\thispagestyle{empty}
\clearpage

\section{The existence question}

A De Bruijn cycle represents every word of a fixed length by an overlapping
cyclic window. Chen, Kitaev, M\"utze and Sun introduced universal partial
words and their cyclic versions, which compress this representation by
allowing wildcards~\cite{CKMS17}. Write $\wild$ for a wildcard, matching either binary
symbol independently at each occurrence. For a cyclic word
$u=u_0\cdots u_{L-1}$ over $\{0,1,\wild\}$, its window at start $i$ is
\[
  W_i=(u_i,u_{i+1},\ldots,u_{i+n-1}),
  \qquad 0\leq i<L,
\]
where subscripts are reduced modulo $L$. A total word $x\in\{0,1\}^n$ is
\emph{compatible} with $W_i$ when $x_k=u_{i+k}$ at every offset $k$ where
the latter character is ordinary. The word $u$ is a universal partial
cycle, or \emph{upcycle}, if every $x$ has exactly one compatible starting
index. Starting indices remain distinct even when windows overlap. The
cycle is nontrivial when it contains both an ordinary symbol and a
wildcard. If each window has $d$ wildcards, its \emph{diamondicity} is $d$.

Fillmore, Goeckner, Kirsch, Martin and McGinnis give binary examples with
$n=8$ and $d=1$ in Theorem 3.1 of \cite{FGKMM26}. Their Question 9.2 asks
whether upcycles with $d>1$ exist. We prove the following.

\begin{theorem}\label{thm:main}
There is a nontrivial binary upcycle for words of length $12$ with
diamondicity $3$ and cyclic length $512$.
\end{theorem}

The example establishes that multiple
independent wildcards in every window can coexist with unique coverage in
a single cycle. The exact object and all verification code appear below.

The closest prior constructions in \cite{FGKMM26} supply the $d=1$ seeds
and multiply alphabet size while preserving $n$ and $d$ (Theorem 5.1 and
Corollary 5.2). The present $d=3$ seed extends the known range of
diamondicities. Its wildcard positions have period four. This meets the
necessary condition $d=n/4$ in their Corollary 8.14 for a nontrivial
period-four \emph{frame}, the pattern recording which positions are wildcards.

Carey, Kearney, Kirsch and Popescu study universal partial quasi-families
in Definition 8.2 of \cite{CKKP25}. Such a family may use several cyclic
words to cover all total words. Here the entire cover is realized by one
cyclic word, including the windows crossing its boundary.

\section{The exact cycle and its proof}

Concatenate the following four rows, in order, to form a string $P$ of
128 octal digits. Row breaks are purely typographical.
\begin{center}
\begin{tabular}{l}
\texttt{00177617563754173657321716771457}\\
\texttt{12371016661655164636441633162656}\\
\texttt{22161556133611160676045602151214}\\
\texttt{46141314063206120412261006540236}
\end{tabular}
\end{center}
For each digit $b$, let $\beta(b)$ be its three-bit binary expansion,
including leading zeroes. Replace $b$ by $\wild\beta(b)$ and concatenate
the resulting blocks to obtain $u$. For example, the digits $0,1,7$
decode as $\wild000,\wild001,\wild111$. The last block is immediately
followed by the first block when $u$ is read cyclically. In the code below,
the ASCII character \texttt{*} represents $\wild$.

For a window $W_i$, let $C_i\subseteq\{0,1\}^{12}$ be its set of
compatible total words. The following elementary criterion reduces unique
coverage to finite checks on ordinary symbols.

\begin{lemma}\label{lem:conflict}
Two partial words of the same length have disjoint sets of binary
completions if and only if at some common offset both characters are
ordinary and unequal.
\end{lemma}
\begin{proof}
An unequal pair of specified bits prevents any simultaneous completion.
Conversely, suppose there is no such pair. At each offset where either
word specifies a bit, choose that bit; the choice is consistent whenever
both specify one. Choose either bit at any remaining offset. This gives a
common completion.
\end{proof}

\begin{proof}[Proof of Theorem~\ref{thm:main}]
The decoding gives $128\cdot4=512$ positions, with wildcards exactly at
indices congruent to zero modulo four. Since four divides 512, this
description is valid across the cyclic boundary. Every 12-position window
therefore contains three wildcards and nine ordinary symbols. The cycle
is nontrivial, and its length exceeds 12; in particular, no physical
position occurs twice within a window. Each $C_i$ has exactly $2^3=8$
members.

For the displayed word, the finite check in the first nested loop of
Listing~\ref{lst:check} verifies
\begin{equation}\label{eq:conflict}
  \forall\,0\leq i<j<512\quad\exists\,0\leq k<12:\quad
  u_{i+k},u_{j+k}\in\{0,1\},\quad u_{i+k}\ne u_{j+k}.
\end{equation}
It checks all $\binom{512}{2}=130816$ pairs, using the characters of the
same decoded word and cyclic subscripts. Every test succeeds. By
Lemma~\ref{lem:conflict}, the sets $C_i$ are pairwise disjoint. Their
union, a subset of $\{0,1\}^{12}$, has cardinality
\[
  \sum_{i=0}^{511}|C_i|=512\cdot8=4096=2^{12}.
\]
Thus these sets partition all binary 12-words. Each total word is
compatible with exactly one starting index, proving the theorem.
\end{proof}

This is a computer-assisted finite proof: the pair check is an exhaustive
integer and character computation, and its full input and implementation
are printed here. The second loop of the same listing directly checks the
conclusion by comparing each of the 4096 total words with all 512 windows.
Both checks use only the displayed word and standard Python operations.

\clearpage
\section{Complete reproducible verification}

Listing~\ref{lst:check} is a complete Python 3 program using only the
standard library. Run it normally, with assertions enabled, for example
as \texttt{python3 verify.py}; do not use the \texttt{-O} option. It
constructs every window with explicit reduction modulo 512, so the eleven
starts $501,\ldots,511$ that wrap around are included. It then checks
both pairwise incompatibility and unique coverage, without sampling.

\lstinputlisting[caption={Complete solver-free verification of the displayed cycle.},
  label={lst:check}]{verify.py}

Successful execution prints
\begin{quote}
\texttt{Verified: 512 windows, 130816 pairs, 4096 words.}
\end{quote}
The direct coverage loop performs $4096\cdot512=2097152$ window
compatibility tests. In particular, it counts starting indices rather
than different window strings or rotation classes. Overlapping windows
read the same characters of $u$; no mismatch predicate is treated as an
independent variable. All counts in this program are exact integers.

\section{Discovery and further questions}

The witness was obtained by an AI-assisted SAT search using exact coverage of
binary 12-words by partial 12-tiles, balance at literal 11-character
overlaps, and guarded cuts to join selected edges into one component.
An Euler circuit then supplied the cyclic word. The finite proof above
verifies the resulting object directly. Separate implementations checked
ordinary conflicts and exhaustive coverage, using published $d=1$ examples
as positive controls and rejecting bit-damaged and all-zero examples.

The theorem supplies an affirmative answer to Question 9.2 of
\cite{FGKMM26}. It also realizes the nondivisibility regime raised in the
paragraph immediately following that question: $12\nmid512=2^{12-3}$.
The next questions concern the possible parameters: does diamondicity
$2$ occur, what is the smallest order supporting diamondicity greater
than one, and can a uniform construction produce unbounded diamondicity?

\section*{AI methodology}

OpenAI Codex provided substantial assistance with the search
formulation and implementation, witness discovery, verification programs,
literature comparison, proof exposition and manuscript preparation.
This assistance materially affected all sections and the code.

\bibliographystyle{alpha}
\bibliography{references}
\end{document}
